\documentclass[11pt,a4paper]{article}
\usepackage[utf8]{vietnam}
\usepackage{amsmath,amssymb}
\usepackage{tabularx}
\usepackage{tikz}
\usepackage{xcolor}
\usepackage{fancyhdr}
\usepackage[top=1.8cm, bottom=2.0cm, left=1.8cm, right=1.8cm]{geometry}

\setlength{\parskip}{3pt}
\setlength{\parindent}{0pt}

\pagestyle{fancy}
\fancyhf{}
\renewcommand{\headrulewidth}{0.4pt}
\lhead{\footnotesize \textbf{CHÙA HƯƠNG SƠN -- HƯỚNG DẪN CHẤM THI}}
\rhead{\footnotesize \textbf{ĐÁP ÁN CHI TIẾT ĐỀ THI THỬ SỐ 02}}
\cfoot{\thepage}

\begin{document}

\begin{center}
{\large\bfseries ĐÁP ÁN \& HƯỚNG DẪN CHẤM ĐỀ THI THỬ CUỐI HỌC KỲ I (ĐỀ SỐ 02)}\\[1.5mm]
{\bfseries ĐƠN VỊ: CHÙA HƯƠNG SƠN \quad|\quad MÔN: TOÁN -- LỚP 9 \quad|\quad NĂM HỌC 2026 -- 2027}\\[1mm]
\textit{(Chuẩn cấu trúc định dạng mới của Sở Giáo dục và Đào tạo)}
\end{center}

\vspace{1mm}
\hrule
\vspace{3mm}

\noindent
\textbf{\large PHẦN I. TRẮC NGHIỆM KHÁCH QUAN (7,0 ĐIỂM)}

\vspace{2mm}
\noindent
\textbf{1. Bảng đáp án Phần 1: Trắc nghiệm nhiều lựa chọn (3,0 điểm)} \textit{(Mỗi câu đúng được 0,25 điểm)}
\begin{center}
\begin{tabular}{|c|c|c|c|c|c|c|c|c|c|c|c|c|}
\hline
\textbf{Câu} & \textbf{1} & \textbf{2} & \textbf{3} & \textbf{4} & \textbf{5} & \textbf{6} & \textbf{7} & \textbf{8} & \textbf{9} & \textbf{10} & \textbf{11} & \textbf{12} \\
\hline
\textbf{Đáp án} & \textbf{A} & \textbf{A} & \textbf{C} & \textbf{A} & \textbf{B} & \textbf{A} & \textbf{B} & \textbf{B} & \textbf{A} & \textbf{B} & \textbf{A} & \textbf{C} \\
\hline
\end{tabular}
\end{center}

\vspace{2mm}
\noindent
\textbf{2. Bảng đáp án Phần 2: Trắc nghiệm Đúng -- Sai (2,0 điểm)}
\begin{center}
\begin{tabular}{|c|c|c|c|c|l|}
\hline
\textbf{Câu} & \textbf{Ý a} & \textbf{Ý b} & \textbf{Ý c} & \textbf{Ý d} & \textbf{Ghi chú đáp án chuẩn} \\
\hline
\textbf{Câu 13} & \textbf{Đ} & \textbf{Đ} & \textbf{S} & \textbf{Đ} & Ý c) sai: Chỗ gặp cách A là $50 \times 2 = 100\text{ km}$ (cách B là $80\text{ km}$). \\
\hline
\textbf{Câu 14} & \textbf{Đ} & \textbf{Đ} & \textbf{Đ} & \textbf{S} & Ý d) sai: Áp dụng HTL $OA^2 = OH \cdot OC \Rightarrow OC = \frac{25}{3}\text{ cm}$. \\
\hline
\end{tabular}
\end{center}

\vspace{2mm}
\noindent
\textbf{3. Bảng đáp án Phần 3: Trắc nghiệm trả lời ngắn (2,0 điểm)} \textit{(Mỗi câu đúng được 0,5 điểm)}
\begin{center}
\begin{tabular}{|c|l|c|c|}
\hline
\textbf{Câu} & \textbf{Nội dung câu hỏi} & \textbf{Kết quả bài làm} & \textbf{Điểm} \\
\hline
\textbf{Câu 15} & Nghiệm bất phương trình $2(x - 3) + 7 < 3(x + 1) - 4$ & \textbf{$x > 2$} & 0,5đ \\
\hline
\textbf{Câu 16} & Giá trị biểu thức $M = \sqrt{(\sqrt{5}-3)^2} + \sqrt{5}$ & \textbf{$3$} & 0,5đ \\
\hline
\textbf{Câu 17} & Khoảng cách từ chân hải đăng đến tàu ($50 / \tan 30^\circ$) & \textbf{$86{,}6\text{ m}$} & 0,5đ \\
\hline
\textbf{Câu 18} & Diện tích bề mặt miếng pizza ($\frac{3{,}14 \times 18^2}{8}$) & \textbf{$127{,}2\text{ cm}^2$} & 0,5đ \\
\hline
\end{tabular}
\end{center}

\vspace{3mm}
\noindent
\textbf{\large PHẦN II. TỰ LUẬN (3,0 ĐIỂM)}

\vspace{2mm}
\textbf{Câu 19 (1,0 điểm).}
\begin{itemize}
  \item[a)] \textit{(0,5 điểm)} Giải hệ: $\begin{cases} 2x - 3y = 1 \\ 3x + 2y = 8 \end{cases} \Leftrightarrow \begin{cases} 4x - 6y = 2 \\ 9x + 6y = 24 \end{cases} \Leftrightarrow \begin{cases} 13x = 26 \\ 2x - 3y = 1 \end{cases} \Leftrightarrow \begin{cases} x = 2 \\ y = 1 \end{cases}$. \hfill \textbf{(0,5đ)}
  \item[b)] \textit{(0,5 điểm)} Quy đồng: $B = \dfrac{x - 4}{\sqrt{x}(\sqrt{x}-2)} : \dfrac{\sqrt{x}+2}{\sqrt{x}} = \dfrac{(\sqrt{x}-2)(\sqrt{x}+2)}{\sqrt{x}(\sqrt{x}-2)} \cdot \dfrac{\sqrt{x}}{\sqrt{x}+2} = 1$. \hfill \textbf{(0,5đ)}
\end{itemize}

\vspace{1mm}
\textbf{Câu 20 (1,5 điểm).}
\begin{itemize}
  \item[a)] \textit{(0,75 điểm)} Vì $OH \perp AB$ tại $H \Rightarrow H$ là trung điểm của $AB \Rightarrow CH$ là đường trung trực của $AB \Rightarrow CA = CB$. Xét $\Delta OAC$ và $\Delta OBC$ có: $OA = OB = R$, $CA = CB$, $OC$ chung $\Rightarrow \Delta OAC = \Delta OBC$ (c.c.c) $\Rightarrow \widehat{OBC} = \widehat{OAC} = 90^\circ \Rightarrow CB \perp OB$ tại $B \in (O) \Rightarrow CB$ là tiếp tuyến của $(O)$. \hfill \textbf{(0,75đ)}
  \item[b)] \textit{(0,75 điểm)} Vì $H$ là trung điểm $AB \Rightarrow AH = 4\text{ cm}$. $\Delta OAH$ vuông tại $H \Rightarrow OH = \sqrt{OA^2 - AH^2} = 3\text{ cm}$. Tam giác $OAC$ vuông tại $A$, đường cao $AH$:
  $OA^2 = OH \cdot OC \Rightarrow OC = \dfrac{OA^2}{OH} = \dfrac{25}{3}\text{ cm}$. Lại có $AC = \sqrt{OC^2 - OA^2} = \dfrac{20}{3}\text{ cm}$. \hfill \textbf{(0,75đ)}
\end{itemize}

\vspace{1mm}
\textbf{Câu 21 (0,5 điểm).}
\begin{itemize}
  \item[a)] \textit{(0,25 điểm)} Cạnh đáy $x > 0$, chiều sâu $h > 0$. $x^2 h = 32 \Rightarrow h = \dfrac{32}{x^2}$. Diện tích lát gạch: $S = x^2 + 4xh = x^2 + \dfrac{128}{x} = x^2 + \dfrac{64}{x} + \dfrac{64}{x} \ge 3\sqrt[3]{x^2 \cdot \dfrac{64}{x} \cdot \dfrac{64}{x}} = 48\text{ m}^2$. Dấu ``='' xảy ra khi $x^2 = \dfrac{64}{x} \Leftrightarrow x = 4\text{ m} \Rightarrow h = 2\text{ m}$. Vậy diện tích nhỏ nhất là $48\text{ m}^2$ khi cạnh đáy $4\text{ m}$, chiều sâu $2\text{ m}$. \hfill \textbf{(0,25đ)}
  \item[b)] \textit{(0,25 điểm)} Ta có $P = \dfrac{1^2}{a} + \dfrac{2^2}{b} + \dfrac{3^2}{c}$. Áp dụng BĐT Cauchy-Schwarz dạng Engel:
  \[
  P \ge \dfrac{(1+2+3)^2}{a+b+c} = \dfrac{36}{6} = 6.
  \]
  Dấu ``='' xảy ra khi $\dfrac{a}{1} = \dfrac{b}{2} = \dfrac{c}{3} = \dfrac{a+b+c}{6} = 1 \Leftrightarrow a = 1, b = 2, c = 3$. Vậy $\min P = 6$. \hfill \textbf{(0,25đ)}
\end{itemize}

\end{document}
